\documentclass[11pt,a4paper]{article}
\usepackage[margin=23mm]{geometry}
\usepackage{amsmath,amssymb,amsthm,microtype}
\usepackage[hidelinks]{hyperref}
\newcommand{\E}{\mathbb E}\newcommand{\Pp}{\mathbb P}
\newtheorem{theorem}{Theorem}\newtheorem{lemma}[theorem]{Lemma}
\title{Two neighboring crossing kernels and an asymmetric larger-noise comparison}
\author{Research note, 28 October 2026}\date{}
\begin{document}\maketitle

The nonnegative kernels in the general crossing identity allow neighboring
counts to be combined before estimating their probabilities. This eliminates
all but one incoming count from the lower bound below. As an application we
prove the central comparison through $\lambda=.36$ when exactly one
conditional law has two-point noise. This is a restricted theorem about
actual iid laws, not a replacement for the general larger-noise comparison.

\section{A two-count reserve}
Use the notation and hypotheses of the retained crossing theorem:
$U_0,U_1$ are centered, independent conditional laws,
\[
 -A_0\le U_0\le a_0,\quad -a_1\le U_1\le A_1,
 \qquad a_0+a_1<1,\quad A_1+a_1\le1.
\]
For mutually independent copies put
\[
 W_j=\sum_{l\le j}U_{1l}+\sum_{l>j}U_{0l},\quad T_j=j+W_j,
 \quad L_j(t)=\E(t-T_j)_+.
\]
Assume $na_0\le t<n(1-a_1)$ and write
\[
 g_i(t)=L_{i-1}(t)-L_i(t)-\Pp(T_i\le t)\ge0.
\]
Nonnegativity is the analytic conclusion of the crossing theorem, not an
arithmetic assumption imposed on test cases.

\begin{theorem}[Paired reserve]\label{thm:pair}
For any integer $1\le k\le n-1$, put $u=t-k$. Then
\begin{align}\label{eq:pair}
 g_k+g_{k+1}
 ={}&u+\Pp(W_k>u)+\E(W_{k-1}-u-1)_+\notag\\
 &-\E\bigl[(u-W_{k+1})\mathbf1_{\{W_{k+1}\le u-1\}}\bigr].
\end{align}
In particular, if $Q$ bounds the last expectation and
$\pi_j=\binom njp^j(1-p)^{n-j}$, then
\begin{equation}\label{eq:weightedpair}
 \sum_{i=1}^n\pi_i g_i\ge
 \min\{\pi_k,\pi_{k+1}\}\{u+\Pp(W_k>u)-Q\}.
\end{equation}
The incoming event in these formulas is closed.
\end{theorem}
\begin{proof}
Add the two local identities to obtain
$g_k+g_{k+1}=L_{k-1}-L_{k+1}-\Pp(T_k\le t)-\Pp(T_{k+1}\le t)$.
Centering and put--call parity give
$L_{k-1}=u+1+\E(W_{k-1}-u-1)_+$.
Also $L_{k+1}+\Pp(T_{k+1}\le t)$ equals the last expectation in
\eqref{eq:pair}, including an atom at $W_{k+1}=u-1$, where its weight is one.
This proves the identity. Drop its nonnegative outgoing stop-loss and then
use $g_i\ge0$ to prove \eqref{eq:weightedpair}. No independence of $W_k$
and $W_{k+1}$ is asserted or needed.
\end{proof}

\begin{lemma}[One hinge bounds the weighted incoming event]\label{lem:hinge}
Let $b=1-u$ and choose $y<b$ with $y\ge-1-u$. For every real $z$,
\[
 (u+z)\mathbf1_{\{z\ge b\}}
       \le\frac{(z-y)_+^2}{(b-y)^2}.
\]
Thus any justified quadratic-hinge bound for $Z=-W_{k+1}$ gives a bound
on $Q$ in Theorem~\ref{thm:pair} without an additional factor.
Use the raw hinge quotient: clipping it at one as for an ordinary probability
is not generally valid for this weighted expectation.
\end{lemma}
\begin{proof}
For $z\ge b$, the derivative of $(u+z)/(z-y)^2$ has the sign of
$-(z+y+2u)$. This is nonpositive because $b+y+2u=1+u+y\ge0$.
The ratio at $b$ is $1/(b-y)^2$. Below $b$ the desired left side is zero.
\end{proof}

There is also a reduction to one conditional sum. With $d=a_0+a_1$ and
$\zeta=1-d$, the auxiliary coupling gives $W_{k+1}\ge W_k-d$ pathwise.
Consequently
\begin{equation}\label{eq:onelaw}
 g_k+g_{k+1}\ge u+\E\left[
 \mathbf1_{\{W_k>u\}}-(u-W_k+d)\mathbf1_{\{W_k\le u-\zeta\}}\right].
\end{equation}
Indeed the incoming event implies $W_k\le u-\zeta$, and its weight is
at most $u-W_k+d$. The affine left tail in \eqref{eq:onelaw} must be kept;
a polynomial for the indicator alone is not a valid minorant of it.

\section{An elementary lower bound for the central binomial masses}
\begin{lemma}\label{lem:mass}
If $n\ge42$, $p\in[2/5,53/100]$ and $k=np$ is an integer, then
\[
 \sqrt n\min\{\pi_k,\pi_{k+1}\}>\frac12,
 \qquad k\ge17,\quad n-k\ge20.
\]
\end{lemma}
\begin{proof}
For any integer-valued variable of variance $V$ and largest atom $M$,
add an independent uniform variable on $[-1/2,1/2]$. The resulting density
is at most $M$ and has variance $V+1/12$. A density bounded by $M$ has
variance at least $1/(12M^2)$: about its mean its tail probability is at
least $(1-2Mr)_+$, and integrate $2r(1-2Mr)_+$ over $r\ge0$.
Thus $M\ge1/\sqrt{12V+1}$.

For the binomial law with integer mean, $\pi_k$ is its largest mass and
$\pi_{k+1}/\pi_k=k/(k+1)\ge17/18$. Since $p(1-p)\le1/4$,
\[
 \sqrt n\min\{\pi_k,\pi_{k+1}\}
 \ge\frac{17}{18\sqrt{3+1/42}}>\frac12.
\]
The last strict inequality is the rational comparison
$289/81>127/42$. The integer bounds follow from $np\ge16.8$
and $n(1-p)\ge19.74$.
\end{proof}

\section{The complete comparison for two asymmetric families}
Write $v=p(1-p)$, $s=(1-p)\E U_0^2+p\E U_1^2$, $\lambda=ns$,
and $Y=B+U$. Put
\[
 \rho=\frac{\E|Y-p|^3}{(v+s)^{3/2}},\qquad
 \rho_B=\frac{1-2v}{\sqrt v}.
\]
At the central integer anchor $t=k=np$ the normal terms agree exactly.
The direct comparison is therefore
\[
 \Delta_C=\sqrt n\{\Pp(Y_1+\cdots+Y_n\le k)-F_B(k)\}
                       +C(\rho_B-\rho).
\]
The signed-curvature lemma from the retained v13 note, whose hypothesis
$t\ge np-(1-p)$ holds here, gives
\begin{equation}\label{eq:master}
 \Delta_C\le-\sqrt n\sum_i\pi_i g_i
             +\frac{\lambda}{4v\sqrt n}+C(\rho_B-\rho).
\end{equation}

\begin{theorem}[Central asymmetric comparison through $\lambda=.36$]\label{thm:family}
Let $n\ge42$, $p\in[2/5,53/100]$, $k=np\in\mathbb Z$, $A=3/5$,
and $0<\lambda\le A^2=9/25$. Consider either of the following centered
conditional laws, with the probabilities uniquely determined by centering:
\begin{enumerate}
\item $U_0\in\{-A,\varepsilon\}$, $U_1=0$,
      where $\varepsilon=\lambda/((n-k)A)$;
\item $U_0=0$, $U_1\in\{-\varepsilon,A\}$,
      where $\varepsilon=\lambda/(kA)$.
\end{enumerate}
For every $C\ge0$ the central comparison is strict. More precisely,
\[
 \Delta_C<-\frac1{240}\quad\hbox{in case 1},\qquad
 \Delta_C<-\frac{25}{48}\lambda\quad\hbox{in case 2}.
\]
Both families satisfy the right-strip bounds $A=.6$, $d\le.072$ and
$H=0$. The theorem includes $\lambda=.36$ and exact atom collisions.
\end{theorem}
\begin{proof}
We first check the cubic term rather than dropping it without a sign.
For either family the inward support is at most $A/17<.072$. More generally,
within the right strip with $d\le.072$, the signs of $Y-p$ on its two
conditional supports are fixed. Expanding the two cubes, using
$\E U_0^3\le a_0\E U_0^2$ and $\E U_1^3\ge-a_1\E U_1^2$, gives
\begin{align*}
 (v+s)^{3/2}(\rho-\rho_B)
 &\ge (3\min(p,1-p)-d)s\notag\\
 &\quad-(1-2v)\left(\frac32s+\frac{3s^2}{8v}\right)\\
 &\ge s\left(\frac{87}{250}-\frac{13}{16}s\right)>0.
\end{align*}
Here $(1+x)^{3/2}\le1+3x/2+3x^2/8$ for $x\ge0$,
$v\ge6/25$ and $s\le(9/25)/42$. Thus $C(\rho_B-\rho)\le0$.

In case 1 put $m=n-k\ge20$ and $\mu=\lambda/A^2\in(0,1]$, so
$\varepsilon=A\mu/m$. If $J$ counts the rare negative jumps in the own
count, then
\[
 W_k=m\varepsilon-(A+\varepsilon)J,
 \quad J\sim\operatorname{Bin}\left(m,\frac\mu{m+\mu}\right).
\]
Consequently $W_k>0$ exactly when $J=0$, and
\begin{equation}\label{eq:ownnegative}
 \Pp(W_k>0)=(1+\mu/m)^{-m}>e^{-\mu}\ge e^{-1}>1/3.
\end{equation}
The last elementary inequality follows from $e<3$.
For the incoming count put $c=A+\varepsilon$, $b=(m-1)\varepsilon$;
then $Z=-W_{k+1}=cJ-b$ with
$J\sim\operatorname{Bin}(m-1,\mu/(m+\mu))$ and
$Q=\E[Z\mathbf1_{\{Z\ge1\}}]$.

If $0<\mu\le2/5$, the incoming event requires $J\ge2$, because
$c-b=A-(m-2)\varepsilon<1$. Hence
\[
 Q\le c\E J(J-1)\le A\mu^2,
 \qquad \Pp(W_k>0)\ge1-\mu.
\]
The latter follows from the union bound on the own rare jumps. Therefore
$\Pp(W_k>0)-Q\ge1-2/5-(3/5)(2/5)^2=63/125>2/15$.

If $2/5\le\mu\le1$, then
\[
 2c-b=A\{2-\mu(1-3/m)\}\le249/250<1,
\]
so the incoming event requires $J\ge3$. For $j\ge3$,
$(cj-b)/(j(j-1))$ decreases with $j$, since its consecutive difference
has the positive numerator $c(j+1)-2b$. Thus
\begin{align}\label{eq:qfactorial}
 Q&\le\frac{3c-b}{6}\E J(J-1)\notag\\
  &=\frac A6\frac{\mu^2\{3-(1-4/m)\mu\}(1-1/m)(1-2/m)}{(1+\mu/m)^2}
  \le\frac A3=\frac15.
\end{align}
For completeness, the $\mu$-dependent fraction increases on $[0,1]$:
its derivative has the positive numerator
$\mu\{6-3(1-4/m)\mu-(1-4/m)\mu^2/m\}$.
At $\mu=1$, the final inequality reduces, with $h=1/m$, to
$(1+2h)(1-h)(1-2h)\le(1+h)^2$; their difference is
$3h+5h^2-4h^3\ge0$. Combining \eqref{eq:ownnegative} and
\eqref{eq:qfactorial} yields the paired reserve $>2/15$.
Theorem~\ref{thm:pair} and Lemma~\ref{lem:mass} give
$\sqrt n\sum_i\pi_i g_i>1/15$.
Since $v\ge6/25$ and $\sqrt n>6$,
\[
 \frac\lambda{4v\sqrt n}<\frac{9/25}{4(6/25)6}=\frac1{16}.
\]
Equation~\eqref{eq:master} proves the first conclusion.

In case 2 put $\mu=\lambda/A^2$, so $\varepsilon=A\mu/k$.
Now $W_k=(A+\varepsilon)J-k\varepsilon$ and $W_k>0$ exactly when
$J\ge1$. Bernoulli's inequality gives
\[
 \Pp(W_k>0)=1-(1+\mu/k)^{-k}
       \ge\frac\mu{1+\mu}\ge\frac\mu2.
\]
Moreover $-W_{k+1}\le(k+1)\varepsilon\le54/85<1$, so $Q=0$.
Therefore the normalized paired loss is greater than $\mu/4$.
Again using \eqref{eq:master},
\[
 \Delta_C<-\frac\lambda{4A^2}+\frac\lambda{4(6/25)6}
           =-\frac{25}{48}\lambda.
\]
All threshold and support assumptions hold: the inward bound is
$\varepsilon\le A/17$, and $t=np$ lies between $na_0$ and $n(1-a_1)$.
\end{proof}

\section{A bounded selected-noise bias}
The following observation is useful independently of the two-point families.
Let $K$ be the unique crossing index from the retained crossing theorem,
and let $\zeta=1-a_0-a_1$. If also
\[
 A_0+a_0<\zeta,\qquad A_1+a_1<\zeta,
\]
then
\begin{equation}\label{eq:bias}
 0\le\E W_{K+1}\le\E(U_0)_+\le a_0,
 \qquad 0\le-\E W_K\le\E(U_1)_-\le a_1.
\end{equation}
To see this, the leave-pair residuals
$R_i=t-(i-1)-\sum_{l<i}U_{1l}-\sum_{l>i}U_{0l}$ decrease by at least
$\zeta$ on every sample. The truncated centered mean
$M_0(x)=\E[U_0\mathbf1_{\{U_0\le x\}}]$ vanishes outside
$[-A_0,a_0)$ and satisfies $0\le-M_0(x)\le\E(U_0)_+$.
At most one residual lies in that interval. The exact bias identity
$\E W_{K+1}=-\sum_i\E M_0(R_i)$ proves the first assertion.
For the other sign, write
$W_K=\sum_i[U_{1i}\mathbf1_{\{T_i\le t\}}+U_{0i}\mathbf1_{\{T_i>t\}}]$.
The second term has mean zero because $U_{0i}$ is independent of $T_i$.
Thus $\E W_K=\sum_i\E M_1(R_i-1)$. Its active interval is
$[1-a_1,1+A_1)$, proving the second assertion. These half-open intervals include
lower-endpoint atoms; the upper endpoint contributes zero by centering.
The acquired right strip has both width bounds, since each width is at
most $.672<.928\le\zeta$.

\section{Scope, dependencies and tests}
The paired identity uses the general crossing theorem in
\texttt{crossing.tex} in the v13 research notes (28 October 2026).
The normalized family theorem additionally uses the signed-curvature lemma
in the adjacent \texttt{counts-tools.tex}; its proof there handles all
counts and finite sample sizes. The present proof of the family theorem
uses no external Poisson comparison and no numerical grid.

The companion exact checks test the identities, the same-sum inequality,
weighted quadratic hinges, the selected-bias caps and actual family
comparisons, including rare jumps and collision endpoints. They preserve
negative controls for a missing closed collision, an inadmissible hinge, and
unjustified clipping of a weighted expectation at one.
The acceptance gates are finite algebraic checks of the inequalities displayed
above, not samples presented as uniform coverage.

No conclusion here removes the general count certificates. The proved
family requires $H=0$, $np$ integral and one noiseless conditional law;
general conditional laws and nonzero phases still need a uniform lower bound
on the paired weighted reserve, or another way to absorb curvature.
\end{document}
